\section{Casos de Estudio (Etapa Digital)}

\begin{tcolorbox}[colback=purple!5!white,colframe=purple!75!black,title=Mapa Mental Rápido (Cópialo)]
\textbf{Zara:} Agilidad extrema, escasez artificial, diseño y producción in-house.\\
\textbf{Barilla:} Efecto Látigo (Bullwhip), variabilidad exagerada de pedidos.\\
\textbf{Hospitales (Virginia Mason):} TPS aplicado a servicios (flujo continuo del paciente).
\end{tcolorbox}

\subsection{Caso Zara (Agilidad en Fast Fashion)}
Este caso evalúa cómo Zara evita los grandes descuentos de la industria de la moda. Anota estos 4 pilares:
\begin{enumerate}
    \item \textbf{Diseño Centralizado:} Decenas de miles de diseños nuevos al año con fuerte feedback diario desde las tiendas (Store Managers).
    \item \textbf{Producción In-House vs Outsourcing:} Las prendas "trendy" (de moda rápida) se fabrican cerca (España/Portugal/Marruecos) para rapidez. Los básicos (ej. poleras blancas) se tercerizan a Asia para costo.
    \item \textbf{Postponement (Pospocisión):} Retrasar el teñido o ensamblaje final hasta conocer qué colores se están vendiendo en tiempo real.
    \item \textbf{Escasez Artificial:} No reponer lo que se agotó. Esto obliga al cliente a comprar inmediatamente ("si no lo llevo ahora, mañana no estará").
\end{enumerate}

\subsection{Caso Barilla (JITD / VMI)}
\begin{itemize}[itemsep=1ex]
    \item \textbf{El Problema:} El \textit{Efecto Látigo} (Bullwhip Effect). Pequeñas variaciones de demanda en tiendas causaban fluctuaciones gigantescas en la fábrica debido a la política de lotes y promociones.
    \item \textbf{La Solución JITD (Just In Time Distribution):} Barilla decide tomar el control del inventario de sus distribuidores. En lugar de recibir "pedidos", Barilla mira el inventario del cliente y le repone solo lo necesario (\textit{Vendor Managed Inventory}).
\end{itemize}

\subsection{Caso Hospital Virginia Mason (TPS en Servicios)}
\begin{itemize}[itemsep=1ex]
    \item \textbf{El Problema:} Médicos perdiendo tiempo buscando insumos, pacientes esperando horas (altos \textit{Mudas}).
    \item \textbf{La Solución (Flujo Continuo):} El paciente no se mueve; los especialistas y equipos orbitan alrededor de él. Uso de "kits" estandarizados de instrumentos. Aplicación directa de TPS a personas.
\end{itemize}
